\section{A uniform rational kernel inequality}
\label{sec:ineq}
\begin{theorem}[Kernel bound for every later truncation]\label{thm:kernel}
For every integer $N\ge2$ and $0\le p,y\le1$,
\begin{equation}\label{eq:kernel-bound}
 0\le K_N(p,y)<\frac98.
\end{equation}
The lower bound is strict when $p>0$ and $y<1$. The strict upper bound includes the continuous boundary values.
\end{theorem}
The proof has two finite base cases and an all-$N$ envelope. We give enough information to reconstruct every rational certificate without a symbolic algebra system.

\subsection{Two explicit polynomial base cases}
Multiplication by the positive denominator shows that $K_N<9/8$ is equivalent to positivity of
\[
 P_N(p,y)=8(1+p)(1+py^2)\{9/8-K_N(p,y)\}.
\]
For $N=2,3$ these are the following polynomials:
\begin{align}
P_2={}&8y^2(1+y)p^3+(8y^3+17y^2+8y+8)p^2\notag\\
&+(9y^2-24y-15)p+9,\label{eq:P2}\\
P_3={}&-8y^2(y+1)(y^2+1)p^4\notag\\
&+(-8y^5-8y^4+16y^3+16y^2-8y-8)p^3\notag\\
&+(24y^3+33y^2+24y+24)p^2
 +(9y^2-32y-23)p+9.\label{eq:P3}
\end{align}
An identity check can avoid rational-function division altogether:
\begin{align}\label{eq:polynomial-check}
(1-y)P_N={}&9(1-y)(1+p)(1+py^2)\notag\\
&-8\{(1-py^2)^N(1+p)-(1-p)^N(1+py^2)\}.
\end{align}
Both verifiers check this identity or independently reconstruct its quotient.

For a polynomial $P(x,y)=\sum c_{kl}x^ky^l$ of bidegree at most $(r,s)$, its tensor Bernstein coefficients on the unit square are
\begin{equation}\label{eq:bernstein}
 B_{ij}=\sum_{k\le i,\,l\le j}c_{kl}
 \frac{\binom ik}{\binom rk}\frac{\binom jl}{\binom sl}.
\end{equation}
The Bernstein basis functions are nonnegative and sum to one. Hence a strictly positive minimum of these coefficients proves strict positivity of $P$ on the whole square, including its boundary.

For a mesh of denominator $d$, first substitute
$p=(u+i)/d$, $y=(v+j)/d$ on each cell. Formula \eqref{eq:bernstein} then gives an exact certificate. The data are summarized by
\begin{center}
\begin{tabular}{lrrrr}
\toprule
Polynomial & Bidegree & Mesh & Coefficients & Smallest coefficient\\
\midrule
$P_2$ & $(3,3)$ & $8\times8$ & $1024$ & $631/12288$\\
$P_3$ & $(4,5)$ & $4\times4$ & $480$ & $11337/40960$\\
\bottomrule
\end{tabular}
\end{center}
Every coefficient is stored as a rational number in the accompanying JSON file. More importantly, the standard-library verifier regenerates all of them from \eqref{eq:P2}--\eqref{eq:bernstein}; it does not trust a stored minimum. These finitely many exact additions and multiplications prove the two base cases.

\subsection{An envelope for binomial differences}
Introduce
\[
 J_m(p,y)=\frac{(1-py^2)^m-(1-p)^m}{1-y},\qquad m\ge2.
\]
For fixed $0<y<1$, differentiation in $p$ gives a unique interior maximum at
\begin{equation}\label{eq:pmax}
 p_m(y)=\frac{1-y^{2/(m-1)}}{1-y^{2m/(m-1)}}.
\end{equation}
Indeed the derivative changes sign exactly when
$1-p=y^{2/(m-1)}(1-py^2)$. Its sign is positive at $p=0$ and negative at $p=1$. Substitution yields the exact envelope
\begin{equation}\label{eq:envelope}
 J_m(p,y)\le B_m(y):=(1+y)
 \left(\frac{1-y^2}{1-y^{2m/(m-1)}}\right)^{m-1}.
\end{equation}
The endpoint values follow by continuity.

\begin{lemma}[The binomial envelope decreases in its index]\label{lem:envelope}
For every $0<y<1$, $B_{m+1}(y)<B_m(y)$ for $m\ge2$.
\end{lemma}
\begin{proof}
Put $k=m-1$, $t=-2\log y>0$ and
$h(v)=\log(1-e^{-(1+v)t})$. Direct differentiation gives $h''(v)<0$ for $v\ge0$. Apart from $\log(1+y)$, the logarithm of \eqref{eq:envelope} is
\[
 k\{h(0)-h(1/k)\}=-\frac{h(1/k)-h(0)}{1/k}.
\]
The secant slope of a strictly concave function decreases as the right endpoint increases. Thus the displayed negative secant slope decreases strictly when $k$ increases. Exponentiation proves the claim.
\end{proof}

\subsection{The fourth envelope is below nine eighths}
Set $y=t^3$, $0\le t\le1$. The inequality $B_4(y)<9/8$ is equivalent away from $t=1$ to
\begin{equation}\label{eq:Qfactor}
9(1-t^8)^3-8(1+t^3)(1-t^6)^3=(1-t^2)^3Q(t)>0,
\end{equation}
where
\begin{align}\label{eq:Q}
Q(t)={}&9t^{18}+27t^{16}-8t^{15}+54t^{14}-24t^{13}+82t^{12}
 -48t^{11}\notag\\
&+84t^{10}-56t^9+60t^8-48t^7+34t^6-24t^5
 +6t^4-8t^3+3t^2+1.
\end{align}
The degree-18 Bernstein coefficients after $t=u/2$ and $t=(1+u)/2$ have respectively the strictly positive minima
\begin{equation}\label{eq:Qmin}
 \frac{192789}{262144},\qquad \frac{690189}{1949696}.
\end{equation}
There are only $38$ coefficients. Formula \eqref{eq:bernstein} with one variable regenerates them. At $y=1$ the envelope limit is
$B_4(1)=2(3/4)^3=27/32<9/8$, so the endpoint is covered as well.

\subsection{The infinite family follows without an infinite computation}
Since $1/(1+v)=\sum_{j\ge0}(1-v)^j/2^{j+1}$ on $[0,1]$,
\begin{equation}\label{eq:mixture}
 K_N(p,y)=\sum_{j=0}^{\infty}2^{-j-1}J_{N+j}(p,y).
\end{equation}
For $y<1$ this follows by subtracting the two convergent geometric series. If $N\ge4$, \eqref{eq:envelope}, Lemma~\ref{lem:envelope} and \eqref{eq:Qfactor} give
\[
 J_{N+j}(p,y)\le B_{N+j}(y)\le B_4(y)<9/8.
\]
Averaging proves $K_N\le B_4<9/8$ for every $N\ge4$. Continuity covers $y=1$. The polynomial base cases cover $N=2,3$, completing Theorem~\ref{thm:kernel}.

Taking expectations in Proposition~\ref{prop:remainder} proves \eqref{eq:nine-eighths}. It also proves $E_N\to g$ throughout the full parameter domain, with a uniform geometric rate. There is no dependence on the masses of positive and negative parts of a signed measure.

\begin{remark}[Nature of the finite computation]
The inequality is a computer-assisted ordinary proof: three stated polynomial identities, a stated rational basis conversion, and $1542$ strictly positive rational numbers close its finite component. The all-index reduction is analytic. The check is reproducible under Python's optimized mode and is not a test of finitely many real sample points. The supplied SymPy replay independently derives the two rational kernels rather than importing their coefficient arrays.
\end{remark}
